\section{Statistics Appendix}
\label{app:stats}

This appendix exposes the confirmatory disclosure the reviews demand: per-condition raw
counts, the actual Wilson intervals, the held-out split, the pre-registered thresholds, a
multiplicity note, a TOST/equivalence treatment for the origin null (A3), and the $K$
normalization detail (B1).

\subsection{A1 --- channel $\to$ cooperation (confirmatory \tC)}
\label{app:a1}

\paragraph{Design and split.} IPD, $n{=}20$ matches per cell, played on the held-out
confirmatory seed offset (100). The analysis plan, the $C>0.8$ lock-in threshold, and the
$0.05$ two-sided $\alpha$ were fixed in \texttt{prereg/a1} before these seeds were run.
Refusals are coded as outcomes; API/transport errors are excluded and tracked as a separate
\texttt{api\_error\_rate}.

\paragraph{Raw counts and Wilson intervals.}
Table~\ref{tab:app-a1} gives lock-in successes $k$ out of $n{=}20$ per cell and the Wilson
95\% interval. The lock-in proportion is $\hat p = k/n$; intervals are the Wilson score
interval \citep{wilson1927probable}, used because the normal approximation degenerates at the
$p{=}0$ cells (L1).

\begin{table}[h]
\centering
\caption{A1 \tC{}: lock-in successes $k$ / $n{=}20$ and Wilson 95\% CI, IPD.}
\label{tab:app-a1}
\small
\begin{tabular}{llcc}
\toprule
rung & payoff & $k/n$ & Wilson 95\% CI \\
\midrule
L0 none      & canonical & 3/20 & [0.05, 0.36] \\
L0 none      & novel     & 3/20 & [0.05, 0.36] \\
L1 signal    & canonical & 0/20 & [0.00, 0.16] \\
L1 signal    & novel     & 0/20 & [0.00, 0.16] \\
L2 free-text & canonical & 12/20 & [0.39, 0.78] \\
L2 free-text & novel     & 8/20 & [0.22, 0.61] \\
L3 private   & canonical & 12/20 & [0.39, 0.78] \\
L3 private   & novel     & 6/20 & [0.15, 0.52] \\
\bottomrule
\end{tabular}
\end{table}

\paragraph{Direct contrast tests (Fisher exact + risk difference).}
We report the headline contrasts directly rather than reading them off interval overlap
(Table~\ref{tab:app-a1-contrast}). The fully significant, robust effect is L2 (free text)
$\gg$ L1 (menu signal): canonical risk difference $+0.60$ $[0.33,0.78]$,
$p{=}3{\times}10^{-5}$. Free text also beats silence on canonical demand (L2 vs L0, $+0.45$,
$p{=}0.008$). \textbf{The ``menu worse than silence'' claim is NOT a confirmatory effect}: the
per-cell L1-vs-L0 contrast ($0/20$ vs $3/20$) is non-significant (Fisher $p{=}0.231$, RD
$-0.15$ $[-0.36,+0.04]$, CI crosses $0$); it reaches $p{=}0.026$ only when pooled across
canonical and novel ($0/40$ vs $6/40$, directional). We therefore lead with L2$\gg$L1 and state
that the menu signal does not improve cooperation.

\begin{table}[h]
\centering
\caption{A1 \tC{}: direct contrast tests (Fisher exact 2-sided, Newcombe risk difference 95\%
CI), $n{=}20$/cell.}
\label{tab:app-a1-contrast}
\small
\begin{tabular}{llcc}
\toprule
demand & contrast & Fisher $p$ & RD [95\% CI] \\
\midrule
canon & L1 vs L0 & 0.231 & $-0.15$ [$-0.36$, $+0.04$] \\
canon & L2 vs L0 & 0.0079 & $+0.45$ [$+0.15$, $+0.66$] \\
canon & L2 vs L1 & $3\!\times\!10^{-5}$ & $+0.60$ [$+0.33$, $+0.78$] \\
novel & L1 vs L0 & 0.231 & $-0.15$ [$-0.36$, $+0.04$] \\
novel & L2 vs L0 & 0.155 & $+0.25$ [$-0.03$, $+0.49$] \\
novel & L2 vs L1 & 0.0033 & $+0.40$ [$+0.16$, $+0.61$] \\
\bottomrule
\end{tabular}
\end{table}

\paragraph{Lock-in threshold sensitivity.}
The L0$\to$L2 cooperation effect is not an artifact of the $C>0.8$ cut. Re-thresholding
match-mean cooperation, the canonical effect is significant at every cut: $C>0.7$, $14/20$ vs
$4/20$, $p{=}0.0036$, RD $+0.50$; $C>0.8$, $12/20$ vs $2/20$, $p{=}0.0022$, RD $+0.50$; $C>0.9$,
$11/20$ vs $2/20$, $p{=}0.0057$, RD $+0.45$. It also holds on continuous mean $C$ ($+0.36$
canonical, $+0.27$ novel). Time-to-first-sustained-cooperation: canonical L2 has $13/20$ matches
ever lock (median first-lock round $0$) versus $5/20$ at L0 (median round $4$)---the channel
raises the lock-in rate and brings it forward. (The match-mean L0 lock-in at $C>0.8$ is $0.10$
here vs.\ $0.15$ in the headline, which uses an end-state metric; the L0$\to$L2 direction is
invariant to that choice.)

\paragraph{Novel-payoff replication.} L2 $8/20$ and L3 $6/20$ both exceed L0 $3/20$ on the
novel payoffs, with L2$\gg$L1 significant ($p{=}0.003$); the effect is in-context reasoning, not
memorized game theory.

\subsection{B1 --- channel $\to$ price collusion (confirmatory \tC)}
\label{app:b1}

\paragraph{Design.} Bertrand duopoly, $n{=}10$--$20$ matches per cell, held-out offset 100;
plan and the competitive/monopoly anchors fixed in \texttt{prereg/b1}. The primary statistic is
the collusion index $K$ (next subsection); per-cell values with bootstrap CIs are in
Table~\ref{tab:b1} of the main text.

\paragraph{Bootstrap 95\% CIs and the channel contrast.}
CIs are percentile bootstrap over match-level $K$ ($5000$ resamples within cell). On canonical
demand the L0/L1 CIs are entirely below $0$ (genuine price war) and the L3 CI is entirely above
$1$ (supra-monopoly: $1.229$ $[0.95,1.54]$); see Table~\ref{tab:b1}. The direct channel contrast
$K(\text{L3}){-}K(\text{L0})$ is $+1.584$ $[1.217,1.974]$ on canonical demand (Mann--Whitney
L3$>$L0 $p{=}4{\times}10^{-7}$) and $+0.274$ $[0.158,0.395]$ on novel demand
($p{=}2{\times}10^{-4}$): the channel raises $K$ on both demand curves, strongly supported.

\paragraph{Generality cells (bootstrap CIs, novel demand, $n{=}12$; L0 vs L3).}
The collusion effect replicates beyond the open spine. GPT-4o: $0.192$ $[0.076,0.310]$ $\to$
$0.810$ $[0.714,0.917]$; GPT-4o$\times$open: $0.244$ $[0.111,0.390]$ $\to$ $0.692$
$[0.570,0.805]$. The Claude-self cell has a \emph{degenerate} CI, $1.286$ $[1.286,1.286]$---zero
variance, $K$ identical in all $12$ matches---which is itself the diagnostic flag of \emph{no}
coordination dynamics (\S\ref{sec:generality}, App.~\ref{app:coord}). Claude$\times$Llama
$0.917$ $[0.832,0.992]\to0.989$ $[0.900,1.076]$; family2-self $0.956$ $[0.829,1.081]$;
family2$\times$open $0.593$ $[0.494,0.698]$.

\paragraph{B1 normalization detail (the $K$ index).}
\label{app:k-norm}
$K = (p - p_{\text{competitive}})/(p_{\text{monopoly}} - p_{\text{competitive}})$, where
$p_{\text{competitive}}$ is the static one-shot Bertrand--Nash price and $p_{\text{monopoly}}$
is the static joint-profit-maximizing price, both computed analytically per demand curve (so
the anchors differ between canonical and novel demand). $K$ is computed per match from the mean
realized price over the scored rounds, then averaged across matches in a cell. \textbf{We do not
clip $K$.} $K>1$ arises because the static monopoly anchor is a \emph{one-shot} benchmark, while
finite-horizon repeated play with \emph{dynamic punishment} (revert-to-price-war threats) can
support prices above it; the $K{=}1.23$ canonical-L3 cell is this supra-monopoly regime. $K<0$
(e.g.\ $-0.64$ at L1) denotes a price war \emph{below} the competitive benchmark. The
canonical-vs-novel gap (e.g.\ $1.23$ vs $0.38$ at L3) quantifies canonical-demand inflation; the
qualitative claim is robust to it.

\subsection{Coordination diagnostic for the no-channel frontier cells (exploratory \tE)}
\label{app:coord}

For the no-channel (L0) Bertrand cells we test whether the two agents' prices \emph{converge /
co-move} over rounds (a coordination signature requiring mutual adjustment) versus being
\emph{independently high from round 1}. With no communication channel, any ``agreement'' could
only be tacit coordination through observed price moves, so the trajectory shape is decisive
(novel Bertrand; $p_{\text{competitive}}{=}10$, $p_{\text{monopoly}}{=}24$, grid $8$--$28$).

\begin{table}[h]
\centering
\caption{\tE{}: no-channel coordination diagnostic. Convergence (shrinking $|p_A-p_B|$ +
positive cross-agent correlation) appears only on the open/GPT spine; the high-$K$ frontier cells
are independent high pricing.}
\label{tab:app-coord}
\small
\begin{tabular}{lccc}
\toprule
cell ($K$) & r1 $|p_A{-}p_B|$ & last & converges? \\
\midrule
Claude self ($1.29$)   & 0.0  & 0.0  & no (flat at 28) \\
Claude$\times$Llama ($0.92$) & 12.5 & 11.7 & no \\
family2 self ($0.96$)  & 5.0  & 10.0 & no \\
GPT-4o self ($0.19$)   & 2.0  & 0.33 & yes (corr $0.59$) \\
open spine ($0.11$)    & 4.3  & 0.18 & yes (corr $0.54$) \\
\bottomrule
\end{tabular}
\end{table}

\noindent\textbf{The Claude-self $K{=}1.29$ cell is independent supracompetitive pricing, not
coordination.} Both Claude agents play price $=28.0$ (grid maximum, \emph{above} the monopoly
price $24$) on every round of every match from round 1; only one price value ($28.0$) appears in
all $404$ agent-rounds; $|p_A-p_B|=0$ throughout, $0/12$ matches show any price change, and the
per-match $K$ slope is $\approx0$ (flat $1.286$ from round 1 to 18). The cross-agent correlation
is undefined (zero variance). This is the \emph{opposite} of a coordination signature: dynamic
collusion looks like the GPT-4o and open-spine rows, where dispersed prices \emph{converge} over
rounds ($|p_A-p_B|$ $4.3\to0.18$, positive correlation $\approx0.5$--$0.6$). We therefore report
the Claude no-channel cell as a fixed high-price \emph{disposition} (a real consumer-harm
\emph{outcome}, $K{=}1.29$) and \emph{not} as tacit coordination. Genuine convergence dynamics
appear only \emph{with} a channel on the open/GPT spine.

\subsection{A3 --- origin moderation: non-significant and under-powered (confirmatory \tC)}
\label{app:a3}

\paragraph{Variance decomposition.} A factorial decomposition of cooperation lock-in yields
$\eta^2_{\text{channel}}=0.54$ and $\eta^2_{\text{origin}}=0.03$---the channel rung explains
$\sim$20$\times$ more variance than the origin grouping.

\paragraph{The moderation null, and why TOST does NOT pass.}
Per the prereg, the powered null is \emph{moderation}: ``origin does not measurably change how
strongly the channel works.'' We operationalize the channel effect per pair-seed as
(mean coop at L2/L3) $-$ (mean coop at L0/L1) and compare same-origin pairs
($\{$same\_origin\_cn, same\_origin\_west, same\_family$\}$) against cross-origin pairs
($\{$cross\_origin, cross\_origin\_2$\}$). The same-origin channel effect is $0.688$ (sd $0.271$,
$n{=}30$); the cross-origin effect is $0.517$ (sd $0.376$, $n{=}20$). The observed moderation
(same $-$ cross) is $+0.171$ with a non-significant Welch test ($p{=}0.090$); the pooled sd is
$0.317$. We then run two one-sided equivalence tests (TOST,
\citealp{lakens2017tost,lakens2018equivalence}):

\begin{table}[h]
\centering
\caption{A3 \tC{}: TOST equivalence for the origin moderation. The test \emph{does not} pass at
any defensible bound; the upper one-sided test fails because $+0.171$ sits near the bound and
$n$ is small.}
\label{tab:app-a3-tost}
\small
\begin{tabular}{lccc}
\toprule
bound $\pm\delta$ & value & $p_{\text{upper}}$ & $p_{\text{TOST}}$ \\
\midrule
$0.5\,$sd (small)  & $\pm0.158$ & 0.551 & \textbf{0.551} (NO) \\
abs $0.20$         & $\pm0.200$ & 0.384 & \textbf{0.384} (NO) \\
abs $0.25$         & $\pm0.250$ & 0.211 & \textbf{0.211} (NO) \\
$0.8\,$sd (medium) & $\pm0.253$ & 0.202 & \textbf{0.202} (NO) \\
\bottomrule
\end{tabular}
\end{table}

\noindent\textbf{We do not claim ``origin doesn't matter'' or ``TOST-equivalent.''} The data
support only (a) \emph{no significant moderation} ($p{=}0.090$) and (b) a large, same-signed
channel effect in both groups ($0.69$ vs.\ $0.52$). The honest claim is a \emph{non-significant,
under-powered moderation}: ``within these four models, in English, with this scaffold, origin is
a weaker driver than the channel, but our $n$ cannot rule out a moderate effect.'' Most of the
cross-origin spread is driven by \texttt{cross\_origin\_2} (a DeepSeek-bearing pair) collapsing
at L2/L3 (lock-in $\approx0.10$) while other pairs reach $\approx0.9$. The exploratory
DeepSeek-partner effect (DeepSeek-containing pairs lock in $\approx0.275$ at L2/L3 vs
$\approx0.85$ for DeepSeek-free pairs) shows the residual variance is \emph{individual-model},
cutting across the origin grouping.

\subsection{Oversight calibration null (exploratory \tE)}
\label{app:calib}

The first-pass non-colluding null was defined from ground-truth $K\le0$ rounds
($n\approx730$ per monitor strength). On these rounds the monitor appears to flag $\approx0.98$
(strong) / $\approx0.88$ (weak); the Pearson correlation between true $K$ and monitor score is
$-0.013$ (strong) and $-0.081$ (weak). This null is contaminated---it is drawn from colluding
runs, because the intended L0 no-channel control was never scored on disk (the monitor reads
only when messages exist)---so these are false-flag rates against a ground-truth-defined null,
\emph{not} a clean operating characteristic.

\paragraph{M1 clean-control correction.} We subsequently built a \emph{designed} message-bearing
but genuinely-non-colluding control and computed a full ROC (\texttt{M1/20260620T143225}).
Against clean negatives the monitor score is discriminative, \textbf{AUC $=0.88$}; the default
operating point sits at FPR $=0.48$, and precision is $0.68$--$0.87$ only at the top
$10$--$25\%$ flag budget. The corrected account is \emph{discriminative-but-mis-thresholded},
superseding the ``flags-everything'' reading.

\paragraph{M2 calibration + governance.} With a calibrated threshold
(\texttt{M2/20260621T012544}, $n{=}48$), the false-positive cost drops from $1.00$ ($27/27$
innocent rounds penalized at the default cut) to $0.00$. Attaching enforcement in this
Bertrand-novel cell left $K$ flat ($\approx0.54$, $\Delta K\approx0$), unlike B5E
($0.33\to0.18$): the deterrent effect of enforcement is cell-dependent. C9 detection
($0.85$--$1.00$) and the B5E enforcement effect are confirmatory~\tC{}; M1/M2 are
exploratory~\tE{}.

\subsection{S2 force/cross and the observability de-confound (exploratory \tE)}
\label{app:s2}

\paragraph{Force/cross.} The honest/omit/lie rates and Wilson intervals are in
Table~\ref{tab:s2} (main text), $n{=}8$--$16$/cell with \texttt{read\_truth}$=1.00$ everywhere
(A always knows B is failing). The raw contrast is a lie rate of $0.75$ $[0.47,0.91]$ when the
ask arrives via the \texttt{none} route (an ambient shared-drive note) versus $0.00$
$[0.00,0.24]$ when it arrives as a direct interpersonal message (\texttt{open}/\texttt{private}),
whose Wilson intervals are disjoint. The \texttt{none} route, however, confounds two factors:
\emph{delivery mode} (ambient artifact vs.\ direct message) and \emph{observability} (whether
the overseer can read the channel). The prior draft attributed the contrast to observability;
the de-confound below shows that attribution does not hold.

\paragraph{Delivery-controlled de-confound (\texttt{S2/20260620T062931}).} We hold the
\emph{delivery} fixed---the misleading ask always arrives as the same ambient shared-drive
note---and manipulate \emph{only} observability/accountability (whether A's handling of the note
is auditable by the overseer). This isolates the variable the main-text table confounded. With
$n{=}30$/cell, the active-lie rates (Wilson 95\% CI) are:

\begin{table}[h]
\centering
\caption{S2 de-confound \tE{}: active-lie successes $k$ / $n{=}30$ and Wilson 95\% CI, delivery
held fixed (ambient drive note).}
\label{tab:app-s2-deconf}
\small
\begin{tabular}{lccc}
\toprule
arm & observability & $k/n$ & Wilson 95\% CI \\
\midrule
threat & unobserved             & 20/30 & [0.49, 0.81] \\
threat & observed / accountable & 15/30 & [0.33, 0.67] \\
benign & unobserved             & 15/30 & [0.33, 0.67] \\
benign & observed / accountable & 18/30 & [0.42, 0.75] \\
\bottomrule
\end{tabular}
\end{table}

\noindent In the threat arm the lie rate moves only from $0.67$ to $0.50$ when handling is made
auditable, with overlapping Wilson intervals (not significant); the benign arm \emph{reverses}
($0.50\to0.60$). Holding delivery fixed, observability/accountability does not significantly
reduce misreporting. We therefore reclassify the Table~\ref{tab:s2} contrast as a
\emph{delivery-mode / request-modality} confound rather than an observability effect: agents
refuse an explicit interpersonal request to lie but passively launder a misleading status
delivered as ambient context, and auditing the channel does not engage the guardrail. This
de-confound was queued as a limitation in the prior draft and is now resolved; the corrected
mechanism is reported in \S\ref{sec:s2}.

\paragraph{Powered de-confound (\texttt{S2/20260621T000133}, $n{=}48$).} To rule out that the
$n{=}30$ threat-arm drift ($0.67\to0.50$) was a real but underpowered effect, we re-ran the
threat / explicit-ask / fixed-delivery cell at $n{=}48$. The active-lie rate is $0.71$
$[0.57,0.82]$ unobserved versus $0.77$ $[0.63,0.87]$ observed/accountable---\emph{flat}, the
observed arm marginally higher. The drift vanishes at higher power, confirming that
observability is not the driver and the corrected request-modality mechanism is robust.

\subsection{Multiplicity note}
\label{app:multiplicity}
The confirmatory spine consists of a small, pre-registered set of primary contrasts (A1: L2 vs
L0 and the L1 reversal; B1: channel vs no-channel on $K$; A3: the origin equivalence test). These
were fixed in advance and are the only claims we label \tC{}; we do not control family-wise error
across the much larger exploratory catalog, all of which is labeled \tE{} and reported as
hypotheses rather than confirmed effects. Where an exploratory contrast is highlighted, its
interval is shown so the reader can judge it directly. We deliberately avoid reporting any
exploratory winner with confirmatory statistics.
